\documentclass[11pt]{article}
\usepackage[margin=25mm]{geometry}
\usepackage{fontspec}
\setmainfont{Latin Modern Roman}
\usepackage{amsmath,amssymb,mathtools}
\usepackage{unicode-math}
\setmathfont{Latin Modern Math}
\usepackage{xurl,hyperref}
\hypersetup{colorlinks=true,urlcolor=blue}
\setcounter{secnumdepth}{0}
\setlength{\parindent}{0pt}
\setlength{\parskip}{5pt}
\setlength{\emergencystretch}{3em}
\newcommand{\sref}[2]{\hyperlink{source:#1:#2}{[S#2]}}
\newcommand{\cataloguescope}[1]{\paragraph{Recorded target.}#1}
\begin{document}
\section{Research Problems in Function Theory — Problem 6.115}\label{statement}

\subsection{Definitions and mathematical statement}

\paragraph{Shared notation.}$\mathbb N=\{1,2,\ldots\}$, and $\mathbb Z,\mathbb Q,\mathbb R,\mathbb C$ denote the integers, rationals, reals and complex numbers. Any other conventions are specified locally.


For a proper plane domain $U$, a Borel set $E\subseteq\partial U$, and $z\in U$, $\omega(z,E,U)$ denotes harmonic measure: the probability that planar Brownian motion begun at $z$ first exits $U$ through $E$. For an arbitrary set, use the associated outer measure; for sets off the boundary, use their intersection with $\partial U$.
A rectifiable curve $\Gamma\subseteq\mathbb C$ has finite length $\mathcal H^1(\Gamma)<\infty$. Let $E\subseteq\Gamma$ satisfy $\mathcal H^1(E)=0$, and let $\Omega\subsetneq\mathbb C$ be any simply connected domain with $z\in\Omega$.
\paragraph{Statement recorded in the supplied source.}
Must $\omega(z,E\cap\partial\Omega,\Omega)=0$ hold? The curve need not coincide with the domain boundary. For $\Omega=\mathbb C$, the boundary is empty and the asserted measure is zero. \sref{statement}{2}
\subsection{Short English statement}
Is every length-zero subset of a rectifiable curve negligible for harmonic measure of every simply connected domain?
\subsection{Sources}
\begin{description}
\item[\hypertarget{source:statement:1}{S1}] ULAM.AI and the UnsolvedMath Contributors, UnsolvedMath: A Curated Collection of Open Mathematics Problems, catalogue entry. CC BY 4.0. \url{https://huggingface.co/datasets/ulamai/UnsolvedMath}.
\item[\hypertarget{source:statement:2}{S2}] W. K. Hayman and E. F. Lingham, Research Problems in Function Theory (2018), Problem 6.115, complete problem and update. \url{https://arxiv.org/abs/1809.07200}.
\end{description}
\label{end:statement}
\end{document}
